\documentclass[11pt]{article}
\usepackage[a4paper,margin=22mm]{geometry}
\usepackage[no-math]{fontspec}
\setmainfont{lmroman10-regular.otf}[BoldFont=lmroman10-bold.otf,ItalicFont=lmroman10-italic.otf,BoldItalicFont=lmroman10-bolditalic.otf]
\usepackage{amsmath,amssymb,amsthm,mathtools,graphicx,xcolor,hyperref,xurl,xspace,ifthen}
\setlength{\emergencystretch}{4em}
\SetMathAlphabet{\mathbf}{normal}{OT1}{cmr}{bx}{n}
\SetMathAlphabet{\mathbf}{bold}{OT1}{cmr}{bx}{n}
\newtheorem{theo}{Theorem}[section]
\newtheorem{lemm}[theo]{Lemma}
\newtheorem{theorem}[theo]{Theorem}
\newtheorem*{remas*}{Remarks}
\providecommand{\og}{«\nobreakspace}
\providecommand{\fg}{\nobreakspace»}
\newtheorem{lem}[theo]{Lemma}
\newtheorem{lemma}[theo]{Lemma}
\newtheorem{prop}[theo]{Proposition}
\newtheorem{coro}[theo]{Corollary}
\newtheorem{corol}[theo]{Corollary}
\newtheorem{conj}[theo]{Conjecture}
\newtheorem{Proposition}[theo]{Proposition}
\newtheorem{theoreme}[theo]{Theorem}
\theoremstyle{definition}
\newtheorem{defi}[theo]{Definition}
\newtheorem{exam}[theo]{Example}
\newtheorem{exem}[theo]{Example}
\newtheorem{rema}[theo]{Remark}
\newtheorem{nota}[theo]{Notation}
\newtheorem*{rema*}{Remark}
\newtheorem*{defi*}{Definition}
\newtheorem*{theo*}{Theorem}
\newtheorem*{lemm*}{Lemma}
\newtheorem*{prop*}{Proposition}
\newcommand{\enoncename}{}
\newtheorem{corpusenonce}[theo]{\enoncename}
\newenvironment{enonce}[2][]{\renewcommand{\enoncename}{#2}\begin{corpusenonce}}{\end{corpusenonce}}
\newenvironment{enonce*}[2][]{\par\medskip\noindent\textbf{#2.}\itshape}{\par\medskip}
\newenvironment{proofc}[1][\proofname]{\begin{proof}[#1]}{\end{proof}}
\newcommand{\Mc}{\mathcal{M}}
\newcommand{\pointir}{.\ ---\ }
\providecommand{\mathof}[1]{\mathrm{#1}}

\numberwithin{equation}{section}
\numberwithin{figure}{section}
\numberwithin{table}{section}
\def\endappendix{}
%~Mouliné par MaN_auto v.0.21.4 2019-04-30 09:02:09


\usepackage[all]{xy}

\hyphenation{auto-morphism
every
Roelcke
peri-od-ic
peri-od-icy
rep-re-sen-ta-tion
}








\newcommand{\A}{\mathcal A}
\newcommand{\N}{\mathbb{N}}
\newcommand{\Z}{\mathbb{Z}}
\newcommand{\Q}{\mathbb{Q}}
\newcommand{\K}{\mathcal{K}}
\newcommand{\CC}{\mathbb{C}}

\newcommand*{\mA}{\mathbf{A}}
\newcommand*{\mB}{\mathbf{B}}
\newcommand*{\mC}{\mathbf{C}}
\newcommand*{\mD}{\mathbf{D}}
\newcommand*{\mE}{\mathbf{E}}
\newcommand*{\mF}{\mathbf{F}}
\newcommand*{\mH}{\mathbf{H}}
\newcommand*{\mZ}{\mathbf{Z}}
%\newcommand*{\mA}{\mathbf{A}}
%\newcommand*{\mA}{\mathbf{A}}

\newcommand{\restrict}[2]{#1 \upharpoonright #2}

\newcommand{\e}{\varepsilon}

\DeclareMathOperator{\Aut}{Aut}
\DeclareMathOperator{\Stab}{Stab}
\DeclareMathOperator{\Age}{Age}

\newcommand{\actson}{\curvearrowright}

\newcommand{\RUCB}{\mathrm{RUC}_{b}}
\newcommand{\ROB}{\mathrm{Ro}_{b}}
\newcommand{\WAP}{\mathrm{WAP}}
\newcommand{\AP}{\mathrm{AP}}

\graphicspath{{./figures/}}

\newcommand*{\mk}{\mkern -1mu}
\newcommand*{\Mk}{\mkern -2mu}
\newcommand*{\mK}{\mkern 1mu}
\newcommand*{\MK}{\mkern 2mu}

\hypersetup{urlcolor=purple, linkcolor=blue, citecolor=red}

\newcommand*{\romanenumi}{\renewcommand*{\theenumi}{\roman{enumi}}}
\newcommand*{\Romanenumi}{\renewcommand*{\theenumi}{\Roman{enumi}}}
\newcommand*{\alphenumi}{\renewcommand*{\theenumi}{\alph{enumi}}}
\newcommand*{\Alphenumi}{\renewcommand*{\theenumi}{\Alph{enumi}}}
\let\oldexists\exists
\renewcommand*{\exists}{\mathrel{\oldexists}}
\let\oldforall\forall
\renewcommand*{\forall}{\mathrel{\oldforall}}

\title{Weak almost periodicity of fully determined finite colorings is equivalent to stability of the embedding pair}
\author{Lionel Nguyen Van Thé}
\date{doi:10.5802/ahl.16; CC BY 4.0}
\begin{document}
\maketitle
\noindent\textit{Editorial note.} This task reuses the independent H2 opinion recorded by root on2026-10-10; this import is not a new mathematical review.\par\medskip
\noindent\textit{Editorial note.} Body477-498 and500-517 were displayed as context in the older sample, but their complete local proofs were explicitly reviewed. Their role is now proof with the original excerpts retained.\par\medskip
\noindent\textit{Editorial note.} The source identifies color values with joint-embedding patterns under fully determined. The inherited reviewer explains the simultaneous group action giving the pattern correspondence; author wording is retained.\par\medskip
\noindent\textit{Editorial note.} Other retained theorem statements supply navigation for author narrative cross-references only. They are not free core results and not separately counted.\par\medskip
\noindent\textit{Editorial note.} ApproxWAP and its unresolved external BYT proof pointer at body713 and later are outside this task and are not declared complete.\par\medskip
\noindent\textit{Editorial scope.} One counted proposition prop:WAPStable, complete in both directions, with the local finite-pattern/Roelcke/coloring bridge. Grothendieck and infinite Ramsey are explicit standard inputs. Later ApproxWAP and the old public fixed-point/Ramsey theorem are outside the counted goal.\par\medskip
\noindent\textit{Allowed standard prerequisites (editorial).}\begin{itemize}
\item Grothendieck double-limit criterion for weak almost periodicity. Body672-675 prints the precise original criterion.. Translate weak compactness into equality of the two iterated limits; this classical theorem is not reproved.
\item Infinite Ramsey theorem for finite colors. Homogenize upper/lower triangular embedding patterns after local Roelcke finiteness is proved.
\item Fraisse ultrahomogeneity and embeddings. Body46 and303-308 retain the author definitions and identification.. Lift the finite embedding witnesses to automorphisms.
\end{itemize}\medskip
On the combinatorial side, the general setting is that of first-order structures in the usual sense of model-theory (see for example~\cite{H} for a standard reference) but for simplicity, we will restrict our attention to the relational setting. Given a first-order relational language (i.e. a family $(R_{i})_{i\in I}$ of symbols together with associated arities $m_{i}\geq 1$), a \emph{structure} $\mA$ is a non-empty set $A$, together with a family of subsets $R^{\mA}_{i}\subset A^{m_{i}}$ for every $i\in I$. Naturally attached to such objects is a notion of isomorphism and of embedding, where an embedding is just an isomorphism onto its image; given two structures $\mA$ and $\mB$, the set of all embeddings of $\mA$ in $\mB$ will be denoted by $\binom{\mB}{\mA}$ (note that this differs from the common notation, according to which $\binom{\mB}{\mA}$ refers to the set of all substructures of $\mB$ isomorphic to $\mA$). A structure is \emph{ultrahomogeneous} when any isomorphism between any two of its finite substructures extends to an automorphism. There is now a rich theory around those objects, starting with the seminal work of Fra\"{\i}ss\'e himself~\cite{Fr0}. For that reason, countable ultrahomogeneous structures are now called \emph{Fra\"{\i}ss\'e} structures (denoted here by $\mF$). In the recent developments of Fra\"{\i}ss\'e theory, a main concern is the study of the interaction between the combinatorics of the set $\Age(\mF)$ of all finite substructures of $\mF$, and the dynamics of the automorphism group $\Aut(\mF)$. The main theorem of~\cite{KPT} is a striking illustration of this.

\begin{theo}\label{thm:fpRP}
Let $\mF$ be a Fra\"{\i}ss\'e structure, $\mathcal X$ be a class of $\Aut(\mF)$-flows such that the class of $\mathcal X$-$\Aut(\mF)$-ambits is closed under suprema and factors, and that every $\Aut(\mF)\actson X\in \mathcal X$ admits some $x\in X$ such that $\Aut(\mF)\actson \overline{\Aut(\mF)\cdot x}\in \mathcal X$. Then $\A:=\{ f\in \RUCB(G) : G\actson \overline{G\bullet f} \in \mathcal X\}$ is a unital, left-invariant, closed C$^{*}$-subalgebra of $\RUCB(\Aut(\mF))$, and TFAE:
\begin{enumerate}
\item Every $\Aut(\mF)$-flow in $\mathcal X$ has a fixed point.
\item For every $\e>0$, $\mF$ has the Ramsey property up to $2\e$ for the finite colorings in $(\A)_{\e}$.
\end{enumerate}

These imply the following equivalent statements:
\begin{enumerate}
\setcounter{enumi}{2}
\item Every zero-dimensional $\Aut(\mF)$-flow in $\mathcal X$ has a fixed point.
\item $\mF$ has the Ramsey property for the finite colorings in $\A$.
\end{enumerate}

When the finite colorings are dense in $\A$, all those statements are equivalent.
\end{theo}

The left side of the above diagrams appears to be particularly well adapted for such an analysis. A \emph{joint embedding} $\langle a,z\rangle$ of two structures $\mA$ and $\mZ$ is a pair $(a,z)$ of embeddings of $\mA$ and $\mZ$ into some common structure $\mC$. Two such pairs $\langle a,z\rangle$ with common range $\mC$ and $\langle a',z'\rangle$ with common range $\mC'$ are \emph{isomorphic} (written $\langle a,z\rangle \cong \langle a',z'\rangle$) when there is an isomorphism $c:\mC\rightarrow \mC'$ so that $a'=c\circ a$ and $z'=c\circ z$. Occasionally, the isomorphism type of a joint embedding $\langle a,z\rangle$ will be referred to as its \emph{joint embedding pattern} and will be written $[a,z]$. In what follows, because the language is assumed to be relational, the joint embeddings which satisfy $C=a(A)\cup z(Z)$ will be the only ones to be considered, without any explicit mention of $\mC$. Note also that the notion of joint embedding $\langle \mA, \mZ^{1},\ldots,\mZ^{k}\rangle$ and joint embedding pattern $[\mA, \mZ^{1},\ldots,\mZ^{k}]$ can be defined in the same way in the case of finitely many structures $\mA$, $\mZ^{1},\ldots,\mZ^{k}$.

\begin{defi}
Let $\K$ be a class of finite structures in some first-order language. It has the \emph{definable Ramsey property} when for every $\mA, \mB\in \K$, every $\mZ\in \K$, there exists $\mC \in \K$ such that for every joint embedding $\langle c,z\rangle$ of $\mC$ and $\mZ$, there is $b\in \binom{\mC}{\mB}$ so that the coloring $a\mapsto [a,z]$ is constant on $\binom{b(\mB)}{\mA}$.
\end{defi}

Note the similarity with the usual Ramsey property. For the combinatorialist, what has just been defined should be thought of as $\mC \rightarrow (\mB)^{\mA}_{\mZ}$. The dynamical meaning of the definable Ramsey property will be made explicit soon, but in view of the fixed point properties described previously, it makes more sense to consider the following weakening first, which will look familiar to the model theorist.

\begin{defi}
Let $\K$ be a class of finite structures in some first-order language, and $\mA, \mZ\in \K$. An \emph{unstable $(\mA,\mZ)$-sequence} is a family of joint embeddings $(\langle a_{m}, z_{n}\rangle)_{m, n\in\N}$ of $\mA$ and $\mZ$ such that there exist two different joint embedding patterns $\tau_{<}$ and $\tau_{>}$ satisfying:
\[
\forall m,n\in \N \qquad (m<n \Rightarrow [a_{m},z_{n}] = \tau_{<})\wedge (m>n \Rightarrow [a_{m},z_{n}] = \tau_{>}).
\]

When there is no unstable $(\mA, \mZ)$-sequence, the pair $(\mA, \mZ)$ is \emph{stable}.
\end{defi}

\begin{defi}
Let $\K$ be a class of finite structures in some first-order language. It has the \emph{stable Ramsey property} when it has the definable Ramsey property restricted to stable pairs. More formally: for every $\mA, \mB\in \K$, every $\mZ^{1},\ldots,\mZ^{k}\in \K$ so that every pair $(\mA, \mZ^{i})$ is stable, there exists $\mC \in \K$ such that for every joint embedding $\langle c,z^{1},\ldots,z^{k}\rangle$, there is $b\in \binom{\mC}{\mB}$ so that for every $i\leq k$, the joint embedding pattern $[a,z^{i}]$ does not depend on $a\in \binom{b(\mB)}{\mA}$.
\end{defi}

\begin{theo}\label{thm:MAPStableRP}
Let $\mF$ be a Fra\"{\i}ss\'e structure with Roelcke-precompact automorphism group. TFAE:
\begin{enumerate}
\item \label{1.6.1} $\Aut(\mF)$ is minimally almost periodic.
\item \label{1.6.2} For every $\mA\in \Age(\mF)$, every $\Aut(\mF)$-invariant equivalence relation on $\binom{\mF}{\mA}$ with finitely many classes is trivial.
\item \label{1.6.3} $\Age(\mF)$ has the stable Ramsey property.
\end{enumerate}
\end{theo}

\begin{theo}\label{thm:fpRoecke}
Let $\mF$ be a Fra\"{\i}ss\'e structure. TFAE:
\begin{enumerate}
\item \label{1.7.1} The flow $\Aut(\mF)\actson R(\Aut(\mF))$ has a fixed point.
\item \label{1.7.2} For every $\mA, \mB, \mZ\in \Age(\mF)$, and every finite coloring $\gamma$ of the joint embedding patterns of $\mA$ and $\mZ$, there exists a joint embedding $\langle b,z\rangle$ such that the coloring $a\mapsto \gamma([a,z])$ is constant on $\binom{b(\mB)}{\mA}$.
\end{enumerate}

When $\Aut(\mF)$ is Roelcke-precompact, these conditions are equivalent to:
\begin{enumerate}
\setcounter{enumi}{2}
\item \label{1.7.3} For every $\mA, \mB, \mZ\in \Age(\mF)$, there exists a joint embedding $\langle b,z\rangle$ such that the coloring $a\mapsto [a,z]$ is constant on $\binom{b(\mB)}{\mA}$.
\end{enumerate}
\end{theo}

\begin{theo}\label{thm:minRoeckeDefRP}
Let $\mF$ be a Fra\"{\i}ss\'e structure with Roelcke-precompact automorphism group. TFAE:
\begin{enumerate}
\item Every minimal subflow of $\Aut(\mF)\actson R(\Aut(\mF))$ is trivial.
\item $\Age(\mF)$ has the definable Ramsey property.
\end{enumerate}
\end{theo}

\section{Fixed points in compactifications and finite~oscillation~stability}
\label{section:oscillationstability}

The left uniformity $\mathcal U_{L}$ is generated by those entourages of the form
\[
V_{U}=\{ (g,h)\in G^{2}: g^{-1}h\in U\},
\]
where $U$ is an open neighborhood of the identity element $e_{G}$. It is induced by any left-invariant metric $d_{L}$ compatible with the topology of $G$ (of course, when $G$ is Polish, there is always such a metric). When $G$ is of the form $\Aut(\mF)$, where $\mF$ is a Fra\"{\i}ss\'e structure whose underlying set is $\N$, a basis of open neighborhoods of $e_{G}$ consists of clopen subgroups of the form $\Stab(\mA)$, where $\Stab(\mA)$ denotes the pointwise stabilizer of $\mA$, and $\mA\subset \mF$ denotes a finite substructure. In this uniformity, two elements $g, h\in G$ are $\mA$-close when $g$ and $h$ agree on $A$. Thus, the corresponding entourage can be seen as a partition of $G$ into sets of the form $g\Stab(\mA)$, and the corresponding quotient space coincides with the usual (algebraic) quotient $G/\Stab(\mA)$.

The right uniformity $\mathcal U_{R}$ is defined in a similar way. It is generated by those sets of the form
\[
V_{U}=\{ (g,h)\in G^{2}: gh^{-1}\in U\},
\]
where $U$ is an open neighborhood of the identity element $e_{G}$. It is induced by any right-invariant metric $d_{R}$ compatible with the topology of $G$. When $G$ is of the form $\Aut(\mF)$ and $\mA$ is a finite substructure of $\mF$, two elements $g, h\in G$ are $\mA$-close when $g^{-1}$ and $h^{-1}$ agree on $A$. The corresponding entourage can be seen as a partition of $G$ into sets of the form $\Stab(\mA)g$, and the corresponding quotient space coincides with the right quotient $\Stab(\mA)\backslash G$. For reasons that will become clear later on, we will see in detail in Section~\ref{section:colorings} how to think about these objects.

The Roelcke uniformity $\mathcal U_{L}\wedge\mathcal U_{R}$ is the finest uniformity that is coarser than the two previous uniformities. When $G$ is of the form $\Aut(\mF)$, a typical uniform neighborhood of this uniformity is indexed by two finite substructures $\mA, \mZ \subset \mF$, and two elements $g,h\in G$ are $(\mA,\mZ)$-close when they are equal in the double quotient $\Stab(\mA)\backslash G / \Stab(\mZ)$. We will see in Section~\ref{section:Roelcke} how to translate this combinatorially.

So far, uniformities were given via a description of their entourages. For those that are induced by compactifications of $G$, another convenient way to produce them is to use algebras of bounded functions. For example, consider the set $\RUCB(G)$ of all bounded uniformly continuous maps from $(G,d_{R})$ to $\CC$ (these maps are called \emph{right-uniformly continuous}). This is a unital C$^{*}$-algebra when equipped with the supremum norm, on which the group $G$ acts continuously by left shift: $g\cdot f(h)=f(g^{-1}h)$. Next, we will follow the terminology from~\cite[Chapter~IV, Section~5]{dV} and will call \emph{left-invariant} the closed C$^{*}$-subalgebras of $\RUCB(G)$ that are invariant under this action.

\begin{theo}\label{thm:fpRPcpct}
Let $\mF$ be a Fra\"{\i}ss\'e structure, and let $\A$ be a unital, left-invariant, closed C$^{*}$-subalgebra of $\RUCB(\Aut(\mF))$. TFAE:
\begin{enumerate}
\item \label{3.1.1} The flow $\Aut(\mF)\actson \Aut(\mF)^{\A}$ has a fixed point.
\item \label{3.1.2} For every $\e>0$, $\mF$ has the Ramsey property up to $2\e$ for the finite colorings in $(\A)_{\e}$.
\end{enumerate}

Those imply the following equivalent statements:
\begin{enumerate}
\setcounter{enumi}{2}
\item \label{3.1.3} Every zero-dimensional factor of $\Aut(\mF)\actson \Aut(\mF)^{\A}$ has a fixed point.
\item \label{3.1.4} $\mF$ has the Ramsey property for the finite colorings in $\A$.
\end{enumerate}

When the finite colorings are dense in $\A$, all those statements are equivalent.
\end{theo}

\subsection{Finite oscillation stability and Ramsey properties}
\label{section:colorings}

Let $\mF$ be a Fra\"{\i}ss\'e structure, whose underlying set is $\N$. As before, for a finite substructure $\mA \subset \mF$, let $\Stab(\mA) \subset \Aut(\mF)$ denote the pointwise stabilizer of $\mA$. Given any $g\in G$, its equivalence class $\bar g$ in the right quotient $\Stab(\mA)\backslash \Aut(\mF)$ is the set of all those elements of $G$ that are $\mA$-close to $g$ (i.e. some sort of ball of radius~$\mA$) relative to the right uniformity, and can be thought of as the restriction $\restrict{g^{-1}}{A}$, an embedding of $\mA$ into $\mF$.

Furthermore, because $\mF$ is ultrahomogeneous, every element of $\binom{\mF}{\mA}$ is of that form. In other words, we can identify $\Stab(\mA)\backslash \Aut(\mF)$ and $\binom{\mF}{\mA}$. In addition, since every element of $\Stab(\mA)\backslash \Aut(\mF)$ can be thought of as a ball for the right uniformity, every coloring of $\binom{\mF}{\mA}$, that is, every map $\bar\chi:\binom{\mF}{\mA}\rightarrow \CC$, can be seen as an element $\chi$ of $\RUCB(\Aut(\mF))$ that is constant on small enough balls and satisfies $\chi(g)=\bar\chi(\bar g)$. Here, we will not usually make any notational distinction between $\chi$ and $\bar\chi$, and by a \emph{coloring in} (resp. \emph{finite coloring in}) $\RUCB(\Aut(\mF))$, we will mean exactly a function $\chi$ of that kind (resp. with finite range). From this point of view, note that even if we allow $\mA$ to range over the set of all finite substructures of $\mF$, every finite set $\mathcal C \subset \RUCB(\Aut(\mF))$ of finite colorings can be seen as a finite set of finite colorings defined on the same set $\binom{\mF}{\mA}$, with values in a common set.

\begin{defi}
Let $\mathcal F \subset \RUCB(\Aut(\mF))$ and $\e>0$. Say that $\mF$ has the \emph{Ramsey property (resp. Ramsey property up to $\e$) for colorings in $\mathcal F$} when for every $\mA, \mB$ in $\Age(\mF)$, and every finite set $\mathcal C \subset \mathcal F$ of finite colorings of $\binom{\mF}{\mA}$, there exists $b\in \binom{\mF}{\mB}$ such that every $\chi\in \mathcal C$ is constant (resp. constant up to $\e$) on $\binom{b(\mB)}{\mA}$.
\end{defi}

\section{Roelcke flows and definable colorings}
\label{section:Roelcke}

\begin{defi}
Let $f : G \rightarrow \CC$. It is \emph{Roelcke} when it is uniformly continuous relative to the Roelcke uniformity on $G$.
\end{defi}

Equivalently, $f$ is Roelcke when it is both right and left uniformly continuous on~$G$. In what follows, we will be particularly interested in \emph{Roelcke-precompact} groups, i.e. groups with compact completion relative to the Roelcke uniformity. In that case, every Roelcke function on $G$ is bounded, and the set $\ROB(G)$ of all Roelcke, bounded, functions is a unital, left-invariant, closed C$^{*}$-subalgebra of $\RUCB(G)$. The corresponding compactification $G^{\ROB(G)}$ will be denoted by $R(G)$. After their introduction in Roelcke and Dierolf~\cite{RD}, Roelcke-precompact groups have shown their utility through the work of Uspenskij~\cite{Us3, Us4, Us5}. More recently, several essential contributions by Tsankov~\cite{Ts}, Ben-Yaacov--Tsankov~\cite{BYT} and Ibarluc\'ia~\cite{I0, I1} have shown that their role is central when studying automorphism groups of Fra\"{\i}ss\'e structures from the model-theoretic point of view. As a matter of fact, Roelcke-precompact groups of the form $\Aut(\mF)$ for $\mF$ Fra\"{\i}ss\'e can be easily characterized combinatorially. Indeed, we have seen in Section~\ref{section:oscillationstability} that a typical entourage of the Roelcke uniformity on $\Aut(\mF)$ is indexed by two finite substructures $\mA$ and $\mZ$ of $\mF$, and that two elements $g,h\in \Aut(\mF)$ are $(\mA, \mZ)$-close when
\[
\Stab(\mA)g\Stab(\mZ)=\Stab(\mA)h\Stab(\mZ).
\]
If we denote by $z$ the identity embedding $\mZ\to \mF$, this means:
\[
\langle \restrict{g^{-1}}{A}, z \rangle\cong \langle \restrict{h^{-1}}{A}, z \rangle.
\]

Here, it will be useful to remember that for a joint embedding $\langle a,z\rangle$, $[a,z]$ refers to its pattern, i.e. its isomorphism type.

\begin{prop}\label{prop:Roelcke-precompact}
Let $\mF$ be a Fra\"{\i}ss\'e structure. Then $\Aut(\mF)$ is Roelcke-precompact iff for every $\mA, \mZ \in \Age(\mF)$, there are only finitely many joint embedding patterns of $\mA$ and $\mZ$.
\end{prop}

\begin{proof}
$\Aut(\mF)$ is Roelcke-precompact iff for every entourage $U$ there are\linebreak $g_{1},\ldots,g_{n}\in \Aut(\mF)$ so that every $g\in \Aut(\mF)$ is $U$-close to some $g_{i}$. From the discussion above, this means that for any two finite substructures $\mA, \mZ$ of $\mF$, $\Aut(\mF)$ can be covered by finitely many $\Stab(\mA)\backslash \Aut(\mF)/\Stab(\mZ)$-classes, which holds exactly when there are only finitely many joint embedding patterns of $\mA$ and $\mZ$.
\end{proof}

\begin{prop}
Let $\mF$ be a Fra\"{\i}ss\'e structure. Then finite colorings are dense in $\ROB(\Aut(\mF))$.
\end{prop}

\begin{proof}
Let $f\in \ROB(\Aut(\mF))$ and fix $\e>0$. Since $f$ is bounded, there is a finite set $Y$ so that the range of $f$ is contained in $(Y)_\e$. By uniform continuity of $f$, there are two finite substructures $\mA, \mZ$ of $\mF$ so that $f$ is constant up to $\e$ on every $\Stab(\mA)\backslash \Aut(\mF)/\Stab(\mZ)$-class. For any such class $P$, choose $ h_P\in P, y_P \in Y$ such that $|y_P - f(h_P)|<\e$. For $g\in G$, set $\bar{\bar g}:= \Stab(\mA)g\Stab(\mZ)$, the equivalence class of $g$ in $\Stab(\mA)\backslash
\Aut(\mF)/\Stab(\mZ)$. Then, the map $\chi : g\mapsto y_{\bar{\bar g}}$ is a finite coloring of $\binom{\mF}{\mA}$. It is in $\ROB(\Aut(\mF))$ because it is constant on the $\Stab(\mA)\backslash \Aut(\mF)/\Stab(\mZ)$-classes, and in addition, for any $g\in G$:
\begin{equation*}
|\chi(g)-f(g)|=|y_{\bar{\bar g}} - f(g)|\leq |y_{\bar{\bar g}}-f(h_{\bar{\bar g}})|+ |f(h_{\bar{\bar g}})-f(g)|<2\e.
\qedhere
\end{equation*}
\end{proof}

Thanks to Theorem~\ref{thm:fpRPcpct}, it follows that under the precompactness assumption of $\Aut(\mF)$, every Roelcke flow has a fixed point iff $\mF$ has the Ramsey property for finite colorings in $\ROB(\Aut(\mF))$. To see how this leads to Theorem~\ref{thm:fpRoecke}, we now turn to a description of those colorings that are in $\ROB(\Aut(\mF))$. In fact, the previous proof already provides such a description. Indeed, if $f$ is assumed to be a finite coloring, then it has to be constant on every $\Stab(\mA)\backslash \Aut(\mF)/\Stab(\mZ)$-class for $\mA, \mZ$ large enough. This means exactly that $f$ can be seen as a finite coloring of the joint embedding patterns of $\mA$ and $\mZ$. Therefore, we have just proved the following.

\begin{prop}\label{prop:Roelcke coloring}
Let $\mF$ be a Fra\"{\i}ss\'e structure, $\mA\in \Age(\mF)$ and $\chi$ be a finite coloring of $\binom{\mF}{\mA}$. Then $\chi\in \ROB(\Aut(\mF))$ iff there is a finite substructure $\mZ$ of $\mF$ such that $\chi(a)$ depends only on $[a,z]$, where $z$ is the identity embedding.
\end{prop}

\subsection{Minimal almost periodicity, weakly almost periodic colorings and~the~stable Ramsey property}
\label{subsection:WAP}

As already mentioned in the introduction, minimal almost periodicity is equivalent to the formally stronger notion of having a fixed point in any distal flow. The corresponding class of ambits is closed under suprema and factors~\cite[Chapter~IV, Section~2.27]{dV}. The corresponding compactification is the so-called \emph{maximal group-like compactification of} $G$~\cite[Chapter~IV, Section~6.18]{dV}, to which is attached the \emph{distal algebra}. Since this algebra contains the almost periodic one, it could have been interesting to use Theorem~\ref{thm:fpRP} to derive a different combinatorial characterization of almost periodicity than the one obtained using the algebra $\mathrm{AP}(G)$. However, we will not do that for two reasons. The first one is that the description of the distal algebra provided by Theorem~\ref{thm:fpRP} does not provide any particularly illuminating way to charactize distal colorings. The second is that an even more general result can be obtained by considering a still larger algebra of functions, namely, the \emph{weakly almost periodic algebra} $\WAP(G)$, consisting of all those $f\in \RUCB(G)$ such that the closure of $G\bullet f$ is weakly compact in the Banach space $\RUCB(G)$. Note that by the following result of Grothendieck (which we only state here for topological groups), this is equivalent to the fact that $G\cdot f$ is weakly compact.

\begin{theo}[{Grothendieck \cite[Proposition~7]{Grothendieck}}]
\label{thm:Gro}
Let $G$ be a topological group and $f\in \RUCB(G)$. Then $f\in \WAP(G)$ iff there are no sequences $(g_{m})_{m\in \N}$, $(h_{n})_{n\in \N}$ of elements of $G$ such that $\displaystyle \lim_{m}\lim_{n}f(g_{m}h_{n})$ and $\displaystyle \lim_{n}\lim_{m}f(g_{m}h_{n})$ both exist and are distinct.
\end{theo}

Applying Theorem~\ref{thm:fpRPcpct}, it follows that when it is Roelcke-precompact, $\Aut(\mF)$ is minimally almost periodic iff $\mF$ has the Ramsey property for finite colorings in $\WAP(\Aut(\mF))$. We now proceed as in Section~\ref{section:Roelcke} to show that this leads to the equivalence $\eqref{1.6.1}\Leftrightarrow \eqref{1.6.3}$ in Theorem~\ref{thm:MAPStableRP}. To do so, we follow the same scheme as for the proof of Theorem~\ref{thm:minRoeckeDefRP}. The first step is to characterize weakly almost periodic colorings combinatorially. Following~\cite{BYT}, this can easily be done thanks to Theorem~\ref{thm:Gro}. Recall first that according to Proposition~\ref{prop:Roelcke coloring}, a finite coloring $\chi$ of $\binom{\mF}{\mA}$ is in $\ROB(\Aut(\mF))$ when there is $\mZ \in \Age(\mF)$ and an embedding $z$ of $\mZ$ in $\mF$ so that $\chi(a)$ depends only on $[a,z]$. We will say then that $\chi$ is \emph{fully determined by $z$} when the converse also holds: $\chi(a)=\chi(a')$ implies $[a,z]=[a',z]$. (In other words, $\chi(a)$ is essentially $[a,z]$.)

\begin{prop}\label{prop:WAPStable}
Let $\mF$ be a Fra\"{\i}ss\'e structure with Roelcke-precompact automorphism group, $\mA$ a finite substructure of $\mF$ and $\chi$ be a finite coloring of $\binom{\mF}{\mA}$. Assume that $\chi\in \ROB(\Aut(\mF))$ is fully determined by $z\in \binom{\mF}{\mZ}$. Then $\chi\in \WAP(\Aut(\mF))$ iff the pair $(\mA, \mZ)$ is stable.
\end{prop}

\begin{proof}
We prove that $\chi\notin \WAP(\Aut(\mF))$ iff the pair $(\mA, \mZ)$ is unstable. If $\chi\notin \WAP(\Aut(\mF))$, consider witness sequences $(g_{m})_{m}, (h_{n})_{n}$ provided by Theorem~\ref{thm:Gro}. For $m\in \N$, define $a_{m}=\restrict{g_{m}^{-1}}{A}$ and $z_{m}=h_{m}\circ z$. Then for $m,n\in \N$, we have $\chi(g_{m}h_{n})=[a_{m},z_{n}]$. By Roelcke-precompactness of $\Aut(\mF)$, this provides a finite coloring of the pairs of naturals, so by the standard Ramsey's theorem, passing to subsequences, we may assume that there are joint embedding patterns $\tau_{<}$ and $\tau_{>}$ so that
\[
\forall m,n\in \N \quad (m<n \Rightarrow [a_{m},z_{n}]=\tau_{<})\wedge (m>n \Rightarrow [a_{m},z_{n}]=\tau_{>}).
\]

In particular, $\lim_{m}\lim_{n}\chi(g_{m}h_{n})=\tau_{<}$, $\lim_{n}\lim_{m}\chi(g_{m}h_{n})=\tau_{>}$ and by choice of $(g_{m})_{m}$ and $(h_{n})_{n}$, $\tau_{<}$ and $\tau_{>}$ are distinct, witnessing that $(\mA, \mZ)$ is unstable.

Conversely, assume that $(\mA, \mZ)$ is unstable, as witnessed by sequences $(a_{m})_{m}$ and $(z_{n})_{n}$, and distinct joint embedding patterns $\tau_{<}$ and $\tau_{>}$. By ultrahomogeneity of $\mF$, we can find, for every $m\in \N$, $g_{m}$ and $h_{m}$ so that $a_{m}=\restrict{g_{m}^{-1}}{A}$ and $z_{m}=h_{m}\circ z$. Then for $m,n\in \N$, we have, as above, $\chi(g_{m}h_{n})=[a_{m},z_{n}]$, so
\[
\lim_{m}\lim_{n}\chi(g_{m}h_{n})=\tau_{<}\neq \tau_{>} = \lim_{n}\lim_{m}\chi(g_{m}h_{n}).
\]
Therefore, $\chi\notin \WAP(\Aut(\mF))$.
\end{proof}

\begin{thebibliography}{MMMMMM}
\bibitem[BYT16]{BYT} Ben Yaacov, Itaï; Tsankov, Todor Weakly almost periodic functions, model-theoretic stability, and minimality of topological groups, Trans. Am. Math. Soc., Volume 368 (2016) no. 11, pp. 8267-8294 \href{https://ahl.centre-mersenne.org/item/AHL_2019__2__149_0/#BYT}{Publisher reference}.
\bibitem[dV93]{dV} de Vries, Jan Elements of topological dynamics, Mathematics and its Applications, 257, Kluwer Academic Publishers, 1993 \href{https://ahl.centre-mersenne.org/item/AHL_2019__2__149_0/#dV}{Publisher reference}.
\bibitem[Fra54]{Fr0} Fraïssé, Roland Sur l’extension aux relations de quelques propriétés des ordres, Ann. Sci. Éc. Norm. Supér., Volume 71 (1954), pp. 363-388 \href{https://ahl.centre-mersenne.org/item/AHL_2019__2__149_0/#Fr0}{Publisher reference}.
\bibitem[Gro52]{Grothendieck} Grothendieck, Alexander Critères de compacité dans les espaces fonctionnels généraux, Am. J. Math., Volume 74 (1952), pp. 168-186 \href{https://ahl.centre-mersenne.org/item/AHL_2019__2__149_0/#Grothendieck}{Publisher reference}.
\bibitem[Hod93]{H} Hodges, Wilfrid Model theory, Encyclopedia of Mathematics and Its Applications, 42, Cambridge University Press, 1993, xiii+772 pages \href{https://ahl.centre-mersenne.org/item/AHL_2019__2__149_0/#H}{Publisher reference}.
\bibitem[Iba16a]{I1} Ibarlucía, Tomas The dynamical hierarchy for Roelcke precompact Polish groups, Isr. J. Math., Volume 215 (2016) no. 2, pp. 965-1009 \href{https://ahl.centre-mersenne.org/item/AHL_2019__2__149_0/#I1}{Publisher reference}.
\bibitem[Iba16b]{I0} Ibarlucía, Tomas Méthodes de théorie des modèles pour l’étude de groupes topologiques (2016) (Ph. D. Thesis) \href{https://ahl.centre-mersenne.org/item/AHL_2019__2__149_0/#I0}{Publisher reference}.
\bibitem[KPT05]{KPT} Kechris, Alexander S.; Pestov, Vladimir G.; Todorcevic, Stevo Fraïssé limits, Ramsey theory, and topological dynamics of automorphism groups, Geom. Funct. Anal., Volume 15 (2005) no. 1, pp. 106-189 \href{https://ahl.centre-mersenne.org/item/AHL_2019__2__149_0/#KPT}{Publisher reference}.
\bibitem[RD81]{RD} Roelcke, Walter; Dierolf, Susanne Uniform structures on topological groups and their quotients, McGraw-Hill International Book Co., 1981, xi+276 pages \href{https://ahl.centre-mersenne.org/item/AHL_2019__2__149_0/#RD}{Publisher reference}.
\bibitem[Tsa12]{Ts} Tsankov, Todor Unitary representations of oligomorphic groups., Geom. Funct. Anal., Volume 22 (2012) no. 2, pp. 528-555 \href{https://ahl.centre-mersenne.org/item/AHL_2019__2__149_0/#Ts}{Publisher reference}.
\bibitem[Usp01]{Us3} Uspenskij, Vladimir V. The Roelcke compactification of groups of homeomorphisms, Topology Appl., Volume 111 (2001) no. 1-2, pp. 195-205 \href{https://ahl.centre-mersenne.org/item/AHL_2019__2__149_0/#Us3}{Publisher reference}.
\bibitem[Usp02]{Us4} Uspenskij, Vladimir V. Compactifications of topological groups, Proceedings of the Ninth Prague Topological Symposium (2001) (2002), pp. 331-346 \href{https://ahl.centre-mersenne.org/item/AHL_2019__2__149_0/#Us4}{Publisher reference}.
\bibitem[Usp08]{Us5} Uspenskij, Vladimir V. On subgroups of minimal topological groups, Topology Appl., Volume 155 (2008) no. 14, pp. 1580-1606 \href{https://ahl.centre-mersenne.org/item/AHL_2019__2__149_0/#Us5}{Publisher reference}.
\end{thebibliography}
\end{document}
